\newpage
\phantomsection
\label{prf:boolean-ring-self-negation}
\LRAProofFor{prop:boolean-ring-self-negation}

\begin{remark*}[Return]
\hyperref[prop:boolean-ring-self-negation]{Return to Theorem}
\end{remark*}

\begin{proposition*}[Self Negation]
In any Boolean ring $R$,
\[
  -p=p \quad \text{for every } p\in R.
\]
\end{proposition*}

\begin{proof}[Professional Standard Proof]
\LRAProofBodyStart
By Proposition~\ref{prop:boolean-ring-self-additive-inverse}, $p+p=0$.
By the additive-inverse clause of R1, $p+(-p)=0$. Both $p$ and $-p$ are
therefore additive inverses of $p$; by uniqueness of additive inverses
(Proposition~\ref{prop:group-inverse-unique}), $-p=p$.
\end{proof}

\begin{proof}[Detailed Learning Proof]
\LRAProofBodyStart
An additive inverse of $p$ is anything that can be added to $p$ to reach
$0$. The ring gives one such element, $-p$. Proposition~\ref{prop:boolean-ring-self-additive-inverse}
hands us a second: $p$ itself, since $p+p=0$. Inverses in an abelian group
are unique, so the two candidates are the same element: $-p=p$. The minus
sign is therefore pure decoration in a Boolean ring.
\end{proof}

\begin{remark*}[Proof structure]
Use characteristic $2$ to make $p$ an additive inverse of itself, then use
uniqueness of additive inverses.
\end{remark*}

\begin{dependencies}
\begin{itemize}
  \item \hyperref[prop:boolean-ring-self-additive-inverse]{Self Additive Inverse / Characteristic 2}.
  \item \hyperref[def:notes-rings-definition-l19]{Ring} axiom R1.
  \item \hyperref[prop:group-inverse-unique]{Uniqueness of additive inverses}.
\end{itemize}
\end{dependencies}

\clearpage
